\documentclass{article}


\usepackage{arxiv}
\usepackage{enumitem}
\usepackage[utf8]{inputenc} % allow utf-8 input
\usepackage[T1]{fontenc}    % use 8-bit T1 fonts
\usepackage{hyperref}       % hyperlinks
\usepackage{url}            % simple URL typesetting
\usepackage{booktabs}       % professional-quality tables
\usepackage{amsfonts}       % blackboard math symbols
\usepackage{nicefrac}       % compact symbols for 1/2, etc.
\usepackage{microtype}      % microtypography
\usepackage{lipsum}
\usepackage{fancyhdr}       % header
\usepackage{graphicx}       % graphics
\usepackage{amsmath}
\usepackage{lyu}
\usepackage{amsthm}
\usepackage{amssymb}
\usepackage{xcolor}
\newtheorem{theorem}{Theorem}[section]
\newtheorem{corollary}{Corollary}[theorem]
\newtheorem{lemma}[theorem]{Lemma}
\theoremstyle{definition}
\newtheorem{definition}{Definition}[section]
\newtheorem{assumption}{Assumption}
\newtheorem{proposition}[theorem]{Proposition}
\newtheorem{remark}[theorem]{Remark}

\pagestyle{fancy}
\thispagestyle{empty}
% \rhead{ \textit{ }} 

% Update your Headers here
% \fancyhead[LO]{}
% 
% \fancyhead[RE]{Firstauthor and Secondauthor} % Firstauthor et al. if more than 2 - must use \documentclass[twoside]{article}


% Update your Headers here
\fancyhead[LO]{Multiscale Nudge}




  
\title{
Multiscale Nudge: Cross-scale Data Assimilation from Macroscopic Observations to Microscopic Dynamics
}

\author{
  Liyao Lyu\\
  Department of Mathematics \\
  University of California, Los Angeles \\
  Los Angeles, CA 90024, USA\\
  \texttt{lyuliyao@math.ucla.edu} \\
  \AND
   Xinyue Yu \\
  Department of Mathematics \\
  University of California, Los Angeles \\
  Los Angeles, CA 90024, USA\\
  \texttt{tracy@math.ucla.edu} \\
  \AND
   Hayden Schaeffer  \\
  Department of Mathematics \\
  University of California, Los Angeles \\
  Los Angeles, CA 90024, USA\\
  \texttt{hayden@math.ucla.edu} \\
  %% Affiliation \\
  %% Address \\
  %% \texttt{email} \\
  %% \And
  %% Coauthor \\
  %% Affiliation \\
  %% Address \\
  %% \texttt{email} \\
  %% \And
  %% Coauthor \\
  %% Affiliation \\
  %% Address \\
  %% \texttt{email} \\
}


\begin{document}
\maketitle
\section{specturm wasserstein flow}
Consider the 
\[
J_{\mathrm{nud}}(\nu,\mu^{\mathrm{obs}}) = \frac{1}{2}\int_{\mathbb{R}^d} \left| \nu(\bz) - \mu^{\mathrm{obs}}(\bz) \right|^2 \intd \bz,
\]
with 
\[
 \frac{\delta J_{\mathrm{nud}}}{\delta \nu}(\bz)  = \nu-\mu^\mathrm{obs}.
\]
Then the associated force field is 
\[
g_\nud = \nabla_\bz (\nu-\mu^\obs).
\]
However, this force field $g$ only determines the descent direction. The actual nudging velocity is further determined by the underlying geometry. Under the classical $W_2$ geometry, the corresponding duality map is defined by
\[
J_\nu(g)
\in
\operatorname*{arg\,min}_{v}
\left\{
\int_{\mathbb{R}^d} g(\bz)\cdot v(\bz)\, d\nu(\bz)
+
\frac12
\int_{\mathbb{R}^d} |v(\bz)|^2\, d\nu(\bz)
\right\}.
\]
The minimizer is given pointwise by
\[
J_\nu(g)(\bz)=-g(\bz).
\]
Hence the induced nudging velocity field is
\[
v_{\mathrm{nud}}(\bz)
=
-g(\bz)
=
-\nabla_{\bz}\bigl(\nu(\bz)-\mu^{\mathrm{obs}}(\bz)\bigr).
\]

Equivalently, from the Jordan--Kinderlehrer--Otto (JKO) perspective, the same $W_2$-gradient flow may be characterized as the minimizing movement scheme
\[
\rho_{k+1}
=
\operatorname*{arg\,min}_{\rho}
\left\{
J_{\mathrm{nud}}(\rho,\mu^{\mathrm{obs}})
+
\frac{1}{2\tau}W_2^2(\rho,\rho_k)
\right\}.
\]
Formally, as $\tau\to 0$, this implicit variational update converges to the continuity equation
\[
\partial_t \rho + \nabla\cdot\bigl(\rho\,v_{\mathrm{nud}}\bigr)=0,
\qquad
v_{\mathrm{nud}}=-\nabla\frac{\delta J_{\mathrm{nud}}}{\delta \rho}.
\]
For the present choice
\[
\frac{\delta J_{\mathrm{nud}}}{\delta \rho}
=
\rho-\mu^{\mathrm{obs}},
\]
this yields
\[
v_{\mathrm{nud}}(\bz)
=
-\nabla_{\bz}\bigl(\rho(\bz)-\mu^{\mathrm{obs}}(\bz)\bigr).
\]
By analogy, we define the spectral Wasserstein cost by
\[
W_\gamma(\nu,\mu)^2
:=
\inf_{\pi\in\Pi(\nu,\mu)}
\gamma\!\left(
\int_{\mathbb R^d\times\mathbb R^d}
(\by-\bx)(\by-\bx)^\top\,d\pi(\bx,\by)
\right),
\]
where $\gamma$ is a norm on the cone $S_d^+$ of positive semidefinite matrices. The corresponding local kinetic cost at $\nu$ is then given by
\[
N_\nu(v)^2
:=
\gamma\!\left(
\int_{\mathbb R^d} v(\bz)v(\bz)^\top\,d\nu(\bz)
\right).
\]
Accordingly, the duality map associated with the spectral geometry is defined by
\[
J_\nu^\gamma(g)
\in
\operatorname*{arg\,min}_{v}
\left\{
\int_{\mathbb R^d} g(\bz)\cdot v(\bz)\,d\nu(\bz)
+
\frac12\,N_\nu(v)^2
\right\}.
\]
If $\gamma(S)=\|S\|_{S_1}$, then, since $S$ is positive semidefinite,
\[
\gamma(S)=\operatorname{tr}(S).
\]
Consequently,
\[
W_\gamma(\nu,\mu)^2
=
\inf_{\pi\in\Pi(\nu,\mu)}
\operatorname{tr}\!\left(
\int (\by-\bx)(\by-\bx)^\top\,d\pi(\bx,\by)
\right)
=
\inf_{\pi\in\Pi(\nu,\mu)}
\int |\by-\bx|^2\,d\pi(\bx,\by)
=
W_2(\nu,\mu)^2.
\]
Thus, the choice $\gamma(S)=\|S\|_{S_1}$ recovers the classical quadratic Wasserstein geometry.
If $\gamma(S)=\|S\|_{S_p}$ for some $1<p<\infty$, then
\[
W_\gamma(\nu,\mu)^2
:=
\inf_{\pi\in\Pi(\nu,\mu)}
\left\|
\int_{\mathbb R^d\times\mathbb R^d}
(\by-\bx)(\by-\bx)^\top\,d\pi(\bx,\by)
\right\|_{S_p}.
\]
This defines an intermediate spectral Wasserstein geometry, interpolating between the classical quadratic Wasserstein distance ($p=1$) and the operator-norm geometry ($p=\infty$). The associated local tangent cost is
\[
N_\nu(v)^2
=
\left\|
\int_{\mathbb R^d} v(\bz)v(\bz)^\top\,d\nu(\bz)
\right\|_{S_p},
\]
which depends on the global covariance structure of the velocity field and is therefore no longer pointwise separable.

The duality map associated with this spectral geometry is defined by
\begin{equation}\label{equ:Lduap_J}
    \begin{aligned}
  J_\nu^{(p)}(g)
\in
\operatorname*{arg\,min}_{v}
\left\{
\int_{\mathbb R^d} g(\bz)\cdot v(\bz)\,d\nu(\bz)
+
\frac12
\left\|
\int_{\mathbb R^d} v(\bz)v(\bz)^\top\,d\nu(\bz)
\right\|_{S_p}
\right\}.  
\end{aligned}
\end{equation}
For $1<p<\infty$, the Schatten-$p$ norm admits the support-function representation
\[
\|S\|_{S_p}
=
\max_{Q\succeq 0,\ \|Q\|_{S_q}\le 1}\operatorname{tr}(QS),
\qquad q=\frac{p}{p-1}.
\]
Accordingly, setting
\[
K_{\gamma_p}
:=
\{Q\succeq 0:\|Q\|_{S_q}\le 1\},
\]
we may rewrite the kinetic term as
\[
\left\|
\int_{\mathbb R^d} v(\bz)v(\bz)^\top\,d\nu(\bz)
\right\|_{S_p}
=
\max_{Q\in K_{\gamma_p}}
\operatorname{tr}\!\left(
Q\int_{\mathbb R^d} v(\bz)v(\bz)^\top\,d\nu(\bz)
\right)
=
\max_{Q\in K_{\gamma_p}}
\int_{\mathbb R^d} v(\bz)^\top Q\,v(\bz)\,d\nu(\bz).
\]
Thus, the Schatten-$p$ geometry may be viewed as selecting, among all admissible quadratic forms $Q\in K_{\gamma_p}$, the one that is most relevant to the current second-order structure of the field.

Now define the force covariance
\[
S_\nu(g)
:=
\int_{\mathbb R^d} g(\bz)g(\bz)^\top\,d\nu(\bz).
\]
If an active matrix $Q_\nu^*$ exists, namely
\[
Q_\nu^*
\in
\operatorname*{arg\,max}_{Q\in K_{\gamma_p}}
\operatorname{tr}\!\bigl(QS_\nu(g)\bigr),
\]
and if $Q_\nu^*$ is invertible, then the duality map admits the representation
\[
J_\nu^{(p)}(g)(\bz)=-(Q_\nu^*)^{-1}g(\bz).
\]
For the endpoint case $p=\infty$, we have
\[
\gamma_\infty(S)=\|S\|_{S_\infty}=\lambda_{\max}(S),
\]
and the corresponding representing set is
\[
K_{\gamma_\infty}
=
\{Q\succeq 0:\|Q\|_{S_1}\le 1\}
=
\{Q\succeq 0:\operatorname{tr}(Q)\le 1\}.
\]
Hence
\[
\|S\|_{S_\infty}
=
\max_{Q\in K_{\gamma_\infty}}\operatorname{tr}(QS).
\]
Approximating \eqref{equ:Lduap_J} by an empirical measure with $n$ particles, the duality-map minimization reduces to the finite-dimensional problem
\[
\min_{V\in\mathbb R^{n\times d}}
\left\{
\frac1n \langle G,V\rangle
+
\frac12
\left\|
\frac1n V^\top V
\right\|_{S_p}
\right\},
\]
where
\[
G=
\begin{bmatrix}
g(\bZ^1)^\to1p\\
\vdots\\
g(\bZ^n)^\top
\end{bmatrix}\in\mathbb R^{n\times d},
\qquad
g(\bZ)=\nabla_{\bz}\frac{\delta J_{\mathrm{nud}}}{\delta \nu}(\bZ),
\]
and $V\in\mathbb R^{n\times d}$ is the particle velocity matrix.
If $G$ admits the singular value decomposition
\[
G=U\operatorname{diag}(\sigma_i)W^\top,
\]
then the minimizer is given by
\[
V=\Xi_p(G)
=
-\|G\|_{S_q}^{\,2-q}\,
U\operatorname{diag}(\sigma_i^{\,q-1})W^\top,
\]
where
\[
r=2p,\qquad q=\frac{r}{r-1}=\frac{2p}{2p-1}.
\]

In particular, if $p=1$, then $r=2$ and $q=2$, so
\[
V
=
-U\operatorname{diag}(\sigma_i)W^\top
=
-G.
\]

If $p=2$, then $r=4$ and $q=\frac43$, and hence
\[
V
=
-\|G\|_{S_{4/3}}^{2/3}\,
U\operatorname{diag}(\sigma_i^{1/3})W^\top.
\]

Finally, if $p=\infty$, then $q=1$, and the formula reduces to
\[
V
=
-\|G\|_{S_1}UW^\top.
\]
Since $\|G\|_{S_1}$ is a scalar, it only sets the overall step size. The update direction is therefore determined by $UW^\top$, which is precisely the Muon-type normalized direction.
\bibliographystyle{unsrt}  

\appendix
\end{document}
